
\subsubsection{6.1 Interpretation of the Directional Result}

The finding that $\tau$\textit{(7.9M) = 20 and $\tau$}(31M) = 50 is consistent with a capacity-diversity trade-off: smaller models may benefit from a more restricted, lower-variance token distribution (achieved through aggressive perplexity filtering) while larger models may leverage distributional diversity from noisier but more varied documents. This is theoretically coherent with Bayesian accounts of in-context learning \citep{Xie2021Explanation}, which predict that n-gram diversity in pre-training data affects ICL capability differently depending on model capacity.

However, caution is warranted in interpreting these results:

\begin{enumerate}
\item The effect size is very small ($\Delta acc$\_norm $\approx$ 0.003--0.006 across the experimental range). At PoC scale (7.9M--18.1M parameters, 500 steps), this may reflect noise from high corpus repetition, short training duration, or random variation.
\end{enumerate}

\begin{enumerate}
\item No formal statistical significance test was conducted. With two seeds per condition, the experiment lacks power for ANOVA. The direction-based gate is a weak criterion.
\end{enumerate}

\begin{enumerate}
\item The repetition factor at $\tau =20$ is extremely high (thousands of passes through ~6,882 documents). Whether the optimal threshold at $\tau =20$ reflects a genuine quality signal or an artifact of fitting to a highly repeated small corpus is unclear.
\end{enumerate}

\begin{enumerate}
\item The 2.$2\times$ scale ratio between 7.9M and 18.1M parameters is modest. The original hypothesis was designed for a 2.$3\times$ ratio at 70M/160M parameters, but the scaling behavior at the PoC scale may not be representative.
\end{enumerate}

\subsubsection{6.2 Failures and Scope Reductions}

The following aspects of the original experimental design were not achieved:

\begin{itemize}
\item \textbf{Scale}: Target 70M/160M parameters; achieved 7.9M/18.1M. Approximately $50\times$ compute gap.
\item \textbf{Token budget}: Target 50B tokens; achieved 1B tokens. $50\times$ gap.
\item \textbf{Corpus}: Target Dolma + FineWeb replication; achieved FineWeb only (Dolma streaming was more complex to implement).
\item \textbf{Statistical test}: Target ANCOVA with p < 0.05 and partial $\eta^2 \geq$ 0.15; achieved direction-based gate only.
\item \textbf{Dedup interaction (P3)}: Not analyzed (dedup had minimal effect in this corpus sample).
\item \textbf{Learning curve analysis (P4)}: Not possible (disk constraints).
\end{itemize}

The preliminary experiment (h-e1) using 7M/16M proxy models with 200 steps produced no signal (ANOVA p=1.0, $\eta^{2}\approx 0)$, establishing that some minimum scale and training duration are necessary for the interaction to be observable.

\subsubsection{6.3 Relation to Prior Work}

The result is consistent with the qualitative finding from Na et al. [2024] that scaling ablation studies across model sizes is necessary to understand data curation effects. It extends beyond existing work by providing a controlled factorial comparison across scales, which prior work (ProX, SoftDedup, REWIRE, WebOrganizer) does not include.

The finding that ''dedup always helps'' --- a widely cited assumption in the field --- could not be evaluated at this scale because FineWeb's sample-10BT is already largely deduplicated (fewer than 100 documents removed per condition). The original hypothesis predicted a sign change in the deduplication effect across scales (P3), which remains untested.

\subsubsection{6.4 Paths Forward}

The full experimental design (70M/160M parameters, 50B tokens, 72 runs, ANCOVA) is fully implemented: the h-e1 pipeline passed all 23 pytest unit tests. The main barrier to the full experiment is compute (~96 A100-hours). At full scale, learning curve dynamics (P4) would also be accessible with appropriate storage.

Post-hoc analysis of the deduplication $\times$ scale interaction from the existing 24-run results.csv is estimated to require approximately one hour and would provide data on P3.
